%% Vol II, Chapter 4 — Platform Support
%% SOURCE: docs/formal/scraps/hardware/platform-integration/PLATFORM_ABSTRACTION.tex
%%         docs/formal/scraps/hardware/platform-integration/L4RE_INTEGRATION.tex
%%         .claude/CLAUDE.md (Platform abstraction, build targets)

\chapter{Platform Support}
\label{vol2:chap:platforms}

StarForth executes across three distinct runtime contexts: the hosted
Linux/POSIX environment, the L4Re/Fiasco.OC microkernel, and bare metal
via LithosAnanke on three hardware architectures (amd64, aarch64, riscv64).
Each context is addressed by a platform abstraction layer that presents a
uniform API to the core VM while confining platform-specific code to
separate compilation units selected at build time.

\section{Platform Abstraction Layer}
\label{vol2:sec:platforms:pal}

%% SOURCE: docs/formal/scraps/hardware/platform-integration/PLATFORM_ABSTRACTION.tex

The platform abstraction layer provides portable timing and clock
functionality across POSIX systems and the L4Re microkernel, selected by a
single Makefile switch. It follows the same vtable pattern as the block I/O
subsystem, achieving zero runtime overhead while keeping platform-specific
code cleanly separated from the core VM.

The core VM is platform-agnostic and speaks only to
\texttt{include/platform\_time.h}. Behind that header, compile-time
selection routes to one of two backends:
\texttt{src/platform/linux/time.c} using \texttt{clock\_gettime}, or
\texttt{src/platform/l4re/time.c} using the RTC server and KIP clock.
Builds choose the path explicitly:

\begin{lstlisting}[language=bash]
make            # POSIX (default)
make L4RE=1     # L4Re/StarshipOS
make MINIMAL=1  # minimal embedded, no platform layer
\end{lstlisting}

\subsection{Platform API}

All timing flows through \texttt{platform\_time.h}. The host calls
\texttt{sf\_time\_init()} once at startup, then reads
\texttt{sf\_monotonic\_ns()} for performance measurement (nanoseconds since
boot) and \texttt{sf\_realtime\_ns()} for wall-clock time (nanoseconds since
the Unix epoch). Additional functions cover clock setting, human-readable
timestamp formatting, RTC availability reporting, and unit conversions
(\texttt{sf\_seconds\_to\_ns}, \texttt{sf\_ns\_to\_ms}, and related
helpers).

Backend selection occurs at compile time on the \texttt{\_\_l4\_\_} define:

\begin{lstlisting}[language=C]
void sf_time_init(void) {
#ifdef __l4__
    sf_time_init_l4re();
    sf_time_backend = &sf_time_backend_l4re;
#else
    sf_time_backend = &sf_time_backend_posix;
#endif
}
\end{lstlisting}

\begin{table}[ht]
  \centering
  \caption{Platform vs.\ feature support}
  \label{tab:platform-features}
  \begin{tabular}{lllll}
    \toprule
    Feature & Linux/POSIX & L4Re/Fiasco.OC & amd64 (LithosAnanke) & aarch64 / riscv64 \\
    \midrule
    Monotonic time   & \texttt{CLOCK\_MONOTONIC} & KIP clock      & TSC + HPET & (see note) \\
    Wall-clock time  & \texttt{CLOCK\_REALTIME}  & RTC server     & APIC timer & (see note) \\
    Interactive REPL & yes                       & yes            & yes        & yes (QEMU validated) \\
    Block I/O        & file + RAM backends       & L4Re dataspace & RAM blocks & RAM blocks \\
    Heartbeat thread & yes                       & yes            & M5 timer   & (see note) \\
    ACL subsystem    & Phases 1--6 complete      & (see note)     & Phase 7 pending & Phase 7 pending \\
    \bottomrule
  \end{tabular}
\end{table}

Entries marked ``(see note)'' are not confirmed in available source material.
%% TODO(bob): fill in aarch64 and riscv64 timing mechanism and heartbeat
%% confirmation once three-arch acceptance logs exist (Phase 7 deliverable).
%% TODO(bob): confirm L4Re ACL subsystem status against StarKernel M7.

\section{Linux / Hosted}
\label{vol2:sec:platforms:linux}

%% SOURCE: src/platform/linux/time.c, .claude/CLAUDE.md

The hosted Linux build is the primary development target. The standard
build procedure produces an optimized binary:

\begin{lstlisting}[language=bash]
make                      # optimized build
make fastest              # ASM + LTO + direct threading
make debug                # symbols enabled
./build/amd64/standard/starforth              # interactive REPL
./build/amd64/standard/starforth --run-tests  # run 936+ tests then REPL
./build/amd64/standard/starforth -c "1 2 + . BYE"  # inline execution
\end{lstlisting}

The POSIX timing backend uses \texttt{CLOCK\_MONOTONIC} for high-precision
monotonic time, \texttt{CLOCK\_REALTIME} for wall-clock time, and
\texttt{localtime}/\texttt{strftime} for timestamp formatting. It always
reports an RTC as available.

Key build flags for the hosted target are documented in
Chapter~\ref{vol2:chap:kernel}. The \texttt{EMERGENCY\_CONSOLE\_ENABLED}
flag (default 1) is relevant to the hosted build as well as the bare-metal
target: setting it to~0 strips the interactive fallthrough surface,
appropriate for production or embedded deployments.

\section{L4Re / Fiasco.OC}
\label{vol2:sec:platforms:l4re}

%% SOURCE: docs/formal/scraps/hardware/platform-integration/L4RE_INTEGRATION.tex

L4Re is a user-level infrastructure layered on the L4 microkernel, providing
memory management, task and thread management, inter-process communication,
device drivers, and runtime libraries. StarForth installs into the L4Re
source tree as \texttt{l4/pkg/starforth}.

\subsection{Package Structure}

The package holds: a \texttt{Control} metadata file (declaring that
the package provides \texttt{starforth} and requires \texttt{libc},
\texttt{libstdc++}, and \texttt{l4re-core}), a top-level \texttt{Makefile},
a \texttt{server/} directory with the L4Re entry point and VM wrapper, a
\texttt{lib/} directory wrapping the existing \texttt{src/} files, the
existing \texttt{include/}, and an \texttt{examples/} client.

\subsection{Memory Management via Dataspaces}

Instead of \texttt{malloc}, the L4Re wrapper backs VM memory with a
dataspace. Initialization allocates a capability slot, allocates a
dataspace of \texttt{VM\_MEMORY\_SIZE}, and maps it into the address space
with \texttt{L4RE\_RM\_F\_SEARCH\_ADDR | L4RE\_RM\_F\_RW}. The mapped
region becomes \texttt{vm\_.memory}. Cleanup detaches the region and frees
the capability.

\subsection{IPC Interface}

The server exposes a small opcode protocol: \texttt{OP\_INTERPRET},
\texttt{OP\_PUSH}, \texttt{OP\_POP}, \texttt{OP\_GET\_STATE}, and
\texttt{OP\_RESET}. \texttt{OP\_INTERPRET} copies an inbound code buffer,
runs \texttt{vm\_interpret}, and returns the VM error status. A matching
client wraps these calls, and registration with the name server makes the
server reachable under the \texttt{starforth} capability.

\subsection{Multi-VM Architecture on L4Re}

A typical deployment runs several FORTH VMs under a root task (Ned), which
acts as name server and VM manager. Distinct VMs serve distinct roles:
one hosting FORTH tasks, one a request/response server, one an interactive
REPL. Ned's \texttt{modules.list} wires each VM's capabilities.

%% TODO(bob): confirm current L4Re port status against StarKernel M7; the
%% L4RE_INTEGRATION scrap describes the L4Re-on-Fiasco path rather than
%% the bare-metal LithosAnanke path, and notes the
%% malloc/free → L4Re-native allocator migration in
%% dictionary_management.c is still pending.

\subsection{L4Re Timing Backend}

The L4Re timing backend draws monotonic time directly from the Kernel
Info Page clock (\texttt{l4\_kip\_clock\_ns(l4re\_kip())}) and real time
from the RTC server offset added to that clock. It obtains the
\texttt{"rtc"} capability from the L4Re environment, queries the offset
over RPC, and degrades gracefully to epoch time when no RTC server is
present. Setting time updates the offset over RPC and requires a write
capability.

\section{Bare Metal: amd64}
\label{vol2:sec:platforms:amd64}

%% SOURCE: .claude/CLAUDE.md (src/starkernel/arch/amd64/)

The amd64 bare-metal path is provided by the LithosAnanke kernel. The
architecture-specific layer resides in \texttt{src/starkernel/arch/amd64/}
and contains:

\begin{itemize}
  \item \texttt{arch.c} --- architecture initialization
  \item \texttt{apic.c} --- Advanced Programmable Interrupt Controller
  \item \texttt{timer.c} --- TSC and HPET timer support
  \item \texttt{interrupts.c} --- interrupt management
  \item \texttt{boot.S} --- low-level boot assembly
  \item \texttt{isr.S} --- interrupt service routine stubs
\end{itemize}

The UEFI loader (\texttt{src/starkernel/boot/uefi\_loader.c}) calls
\texttt{ExitBootServices()} and transfers control to
\texttt{src/starkernel/kernel\_main.c}. Milestone M5 establishes the
100~Hz heartbeat timer using TSC, HPET, and the APIC timer, which drives
the adaptive heartbeat subsystem described in Volume~I.

Build and QEMU execution:

\begin{lstlisting}[language=bash]
make -f Makefile.starkernel                        # build UEFI image
make -f Makefile.starkernel qemu                   # run under QEMU/OVMF
make -f Makefile.starkernel STARFORTH_ENABLE_VM=1  # enable M7 VM integration
\end{lstlisting}

Output: \texttt{build/amd64/kernel/starkernel\_loader.efi} (UEFI PE32+
executable).

\section{Bare Metal: aarch64 and riscv64}
\label{vol2:sec:platforms:arm-riscv}

%% SOURCE: .claude/CLAUDE.md (Build targets section)

Cross-compilation targets for aarch64 and riscv64 are provided via
\texttt{Makefile.starkernel}:

\begin{lstlisting}[language=bash]
make -f Makefile.starkernel ARCH=aarch64    # aarch64 build
make -f Makefile.starkernel ARCH=riscv64    # riscv64 build
make rpi4-cross                             # Raspberry Pi 4 cross-compile
\end{lstlisting}

QEMU-based validation is used for all three architectures. Milestones M0
through M5 have been verified under QEMU/OVMF across all three. The
\texttt{QEMU\_BASELINE.log} file captures reference output for regression
detection.

%% TODO(bob): describe aarch64-specific bring-up notes (Raspberry Pi 4)
%% and riscv64 details once platform-specific documentation is available.
%% Phase 7 ACL acceptance logs will cover three-arch validation.

\section{Raspberry Pi Build}
\label{vol2:sec:platforms:rpi}

%% SOURCE: .claude/CLAUDE.md (Build commands section)

\begin{lstlisting}[language=bash]
make rpi4-cross    # Cross-compile for Raspberry Pi 4 (aarch64)
\end{lstlisting}

%% TODO(bob): describe Raspberry Pi 4 board bring-up details, UART
%% configuration, and any aarch64-specific optimizations once
%% ARM64_OPTIMIZATIONS documentation is promoted to formal tier.

\section{Future Platforms}
\label{vol2:sec:platforms:future}

%% SOURCE: docs/formal/scraps/hardware/platform-integration/PLATFORM_ABSTRACTION.tex

The platform abstraction design anticipates additional backends. Adding a
new backend requires: a new \texttt{time\_<platform>.c} translation unit,
its backend vtable entry, a detection branch in \texttt{platform\_init.c},
and a Makefile flag. Anticipated future targets noted in the platform
abstraction design include bare-metal hardware timers (beyond the current
amd64/APIC path), WebAssembly via \texttt{performance.now()}, RTOSes such
as FreeRTOS and Zephyr, and additional microkernels such as seL4.
